\section{How This Role Fits the Overall Picture}

\cvitem{Test Process}{
\textbf{Roll: Test Architect / Test Lead.} \\
\textbf{Ansvar:} äga teststrategin över samtliga nivåer -- unit, komponent, integration, systemfunktionstest
(SFT) och e2e -- och avgöra vad som måste vara mänskligt granskat kontra vad AI-genererade tester täcker
tillräckligt bra. \\
\textbf{Process:} definierar teststrategi och riskområden innan kod skrivs, kör/granskar testerna löpande
under AI-assisterad utveckling, och stoppar regressioner innan de når en människa. \\
\textbf{Grundat i:} ISTQB Certified Tester, Foundation Level (2006); Test Leader, J2EE-projekt (Sony
Ericsson); on-site systemintegrationstestare i Norge på ett BPM-system som senare tog silver i BPM
Global Awards.
}

\cvitem{Development Process}{
\textbf{Roll: Senior Developer / Code Reviewer.} \\
\textbf{Ansvar:} kodkvalitet, läsbarhet och att AI-genererad kod faktiskt passar in i befintlig
arkitektur -- inte bara att den kör. \\
\textbf{Process:} granskar och styr AI-genererade ändringar med samma rigor som en junior utvecklares
pull request, iterativt under hela utvecklingsflödet. \\
\textbf{Grundat i:} 20+ års praktisk utveckling -- Java/JEE, .NET, React/Angular/Vue, fullstack -- hos
Telenor, IKEA IT, Cybercom, Bouvet Syd, Verisure, ADBSafegate och Zaplox.
}

\cvitem{Application Management Process}{
\textbf{Roll: Application Owner / Förvaltningsledare.} \\
\textbf{Ansvar:} förstå och äga konsekvenserna av ändringar i ett system som redan kör i
produktion -- inte bara leverera en sprint. \\
\textbf{Process:} spårar vad som ändras, varför, och vad det påverkar nedströms, även när
ändringarna är AI-assisterade. \\
\textbf{Grundat i:} Java Developer \& Coordinator, IKEA IT (bryggan mellan central IT och butiksdrift);
GDPR-koordinator, Verisure (ägde dataflödes-analysen i en live-applikation).
}

\cvitem{Version Control \& Configuration Management (CM)}{
\textbf{Roll: Release Manager / CM-ägare.} \\
\textbf{Ansvar:} branch-disciplin, commit-hygien och att ändringar går att rulla tillbaka -- så
att tempo aldrig sker på bekostnad av en ren, spårbar historik. \\
\textbf{Process:} definierar branch-strategi och release-flöde, granskar att AI-genererade commits följer
dem innan merge. \\
\textbf{Grundat i:} DevOps Lead \& Release Manager, IKANO Bank (CA Lisa/Nolio release-automation);
Git/Gerrit, GitLab CI, Jenkins, SonarQube och Artifactory genom hela karriären.
}

\cvitem{Architecture}{
\textbf{Roll: Solution Architect.} \\
\textbf{Ansvar:} de strukturella besluten -- vad en AI-assistent \emph{inte} ska få avgöra själv --
innan kod genereras inom den strukturen. \\
\textbf{Process:} definierar målarkitektur och gränssnitt först, låter sedan AI-verktyget
implementera inom de ramarna, med arkitekturgranskning som ett eget, explicit steg. \\
\textbf{Grundat i:} Solution Architect, IKEA IT (SaaS-utvärdering till Enterprise Architecture Board);
mall för Software Architectural Description (TOGAF, Zachman, Kruchten 4+1) hos Telenor; Application Lead
för ett EAI-integrationslager med 15 system (Hi3G/Tre).
}

\cvitem{DevOps \& Process Coaching}{
\textbf{Roll: DevOps Coach / Scrum Master.} \\
\textbf{Ansvar:} göra AI till en verklig del av teamets arbetssätt -- process, verktyg och
kvalitetsgrindar inräknat -- inte en engångspryl. \\
\textbf{Process:} facilitera samma typ av retrospektiv, ways-of-working-justeringar och coachning som
tidigare, nu applicerat på AI-assisterat arbete. \\
\textbf{Grundat i:} DevOps Coach \& Scrum Master, Nordea; Problem Manager (ITIL), Länsstyrelsen;
DevOps/CD-verktygsutrullning, IKANO Bank.
}
